\section{Tirage \emph{sans remise}}
\label{sec:tirage-sans-remise}

\subsection{Idée générale}
On veut tirer des éléments sans jamais retomber deux fois sur le même
(par exemple, des cartes d'un paquet). Le principe le plus simple et
robuste consiste à :
\begin{enumerate}
  \item Stocker tous les éléments dans un tableau (paquet).
  \item À chaque tirage, tirer un index aléatoire parmi les positions encore disponibles $[0\,..\,\texttt{borne}]$.
  \item Échanger l'élément tiré avec celui en fin de zone disponible (\texttt{borne}), puis réduire la borne d'1.
\end{enumerate}
Ainsi, la zone \texttt{[0..borne]} contient encore les éléments
disponibles, et la zone \texttt{[borne+1..fin]} contient les éléments
déjà tirés.

\subsection{Schéma pas à pas}
Exemple avec 6 cartes initiales : \texttt{[A, B, C, D, E, F]} (indices \texttt{0..5}), \texttt{borne = 5}.

\paragraph{Tirage \#1} on tire un index aléatoire \texttt{i = 2} (la carte \texttt{C}) dans \texttt{[0..5]}.
\[
\begin{array}{c|c|c|c|c|c}
\texttt{0} & \texttt{1} & \color{blue}{\texttt{2}} & \texttt{3} & \texttt{4} & \texttt{5}\\
\boxed{\texttt{A}} & \boxed{\texttt{B}} & \boxed{\texttt{C}} & \boxed{\texttt{D}} & \boxed{\texttt{E}} & \boxed{\texttt{F}} \\
\hline
\end{array}
\]
On échange la case \texttt{i=2} avec la case \texttt{borne=5} :
\[
\begin{array}{c|c|c|c|c|c}
\texttt{0} & \texttt{1} & \color{blue}{\texttt{2}} & \texttt{3} & \texttt{4} & \color{gray}{\texttt{5}}\\
\boxed{\texttt{A}} & \boxed{\texttt{B}} & \boxed{\texttt{F}} & \boxed{\texttt{D}} & \boxed{\texttt{E}} & \fbox{\texttt{C}} \\
\hline
\end{array}
\]
Puis \texttt{borne ← 4}. La carte \texttt{C} est sortie (à droite), elle ne sera plus retirable.

\paragraph{Tirage \#2} on tire \texttt{i = 1} dans \texttt{[0..4]} (car \texttt{borne=4}) :
\[
\begin{array}{c|c|c|c|c|c}
\texttt{0} & \color{blue}{\texttt{1}} & \texttt{2} & \texttt{3} & \color{gray}{\texttt{4}} & \color{gray}{\texttt{5}}\\
\boxed{\texttt{A}} & \boxed{\color{blue}{\texttt{B}}} & \boxed{\texttt{F}} & \boxed{\texttt{D}} & \boxed{\texttt{E}} & \fbox{\texttt{C}} \\
\hline
\end{array}
\]
On échange avec \texttt{borne=4}, puis \texttt{borne ← 3} :
\[
\begin{array}{c|c|c|c|c|c}
\texttt{0} & \texttt{1} & \texttt{2} & \texttt{3} & \color{gray}{\texttt{4 (tirée)}} & \color{gray}{\texttt{5 (tirée)}}\\
\boxed{\texttt{A}} & \boxed{\texttt{E}} & \boxed{\texttt{F}} & \boxed{\texttt{D}} & \fbox{\texttt{B}} & \fbox{\texttt{C}} \\
\hline
\end{array}
\]
On poursuit ainsi jusqu'à ce que \texttt{borne < 0} (tous les éléments tirés), ou qu'on ait tiré le nombre voulu d'éléments.

\subsection{Erreurs fréquentes à éviter}
\begin{itemize}
  \item Ne pas réduire la borne : risque de re-tirer le même élément.
  \item Tirer sur toute la longueur \texttt{[0..N-1]} à chaque fois : doublons possibles.
\end{itemize}

\section{Temporisation avec \texttt{sleep}}

La fonction \texttt{sleep()}, déclarée dans \texttt{<unistd.h>}, suspend
l'exécution du programme pendant un nombre de secondes donné. Elle est
utile pour marquer une pause entre deux affichages (par exemple, laisser
le temps de lire un message avant d'effacer l'écran).

\begin{lstlisting}[language=C]
#include <unistd.h>

printf("Préparez-vous...\n");
sleep(3); // attente bloquante de 3 secondes
printf("C'est parti !\n");
\end{lstlisting}

\begin{UPSTIinfor}{Pour une attente plus courte}
La fonction \texttt{usleep()}, également dans \texttt{<unistd.h>}, attend
en microsecondes (\texttt{usleep(500000)} attend 0,5 seconde).
\end{UPSTIinfor}

\section{Couleurs dans la console}

Un terminal qui interprète les séquences ANSI accepte des codes
d'échappement pour changer la couleur du texte ou du fond avant un
affichage, et pour revenir à la normale ensuite. Le code général est
\texttt{"\textbackslash033[<code>m"}, suivi en fin d'affichage de
\texttt{"\textbackslash033[0m"} pour réinitialiser.

\begin{lstlisting}[language=C]
printf("\033[31mTexte rouge\033[0m\n");   // texte rouge
printf("\033[42m  \033[0m\n");            // carré de fond vert
\end{lstlisting}

Quelques codes de fond courants : \texttt{40} noir, \texttt{41} rouge,
\texttt{42} vert, \texttt{44} bleu. Les mêmes valeurs diminuées de 10
(\texttt{30}, \texttt{31}, \texttt{32}, \texttt{34}) s'appliquent à la
couleur du texte plutôt qu'à son fond.

\begin{UPSTIinfor}{Effacer le terminal}
L'appel \texttt{system("clear");} efface le contenu affiché dans le
terminal, utile en complément d'une temporisation pour rafraîchir
l'affichage d'un plateau de jeu.
\end{UPSTIinfor}
